\documentclass[10pt, letterpaper]{article}

% Packages:
\usepackage[
    ignoreheadfoot,
    top=2 cm,
    bottom=2 cm,
    left=1.2 cm,
    right=1.2 cm,
    footskip=1.0 cm
]{geometry}
\usepackage{titlesec}
\usepackage{tabularx}
\usepackage{array}
\usepackage[dvipsnames]{xcolor}
\definecolor{primaryColor}{RGB}{0, 0, 0}
\usepackage{enumitem}
\usepackage{fontawesome5}
\usepackage{amsmath}
\usepackage[
    pdftitle={Reece Palmer's CV},
    pdfauthor={Reece Palmer},
    pdfcreator={LaTeX with RenderCV},
    colorlinks=true,
    urlcolor=primaryColor
]{hyperref}
\usepackage[pscoord]{eso-pic}
\usepackage{calc}
\usepackage{bookmark}
\usepackage{lastpage}
\usepackage{changepage}
\usepackage{paracol}
\usepackage{ifthen}
\usepackage{needspace}
\usepackage{iftex}

% Ensure that generated PDF is machine readable / ATS parsable:
\ifPDFTeX
    \input{glyphtounicode}
    \pdfgentounicode=1
    \usepackage[T1]{fontenc}
    \usepackage[utf8]{inputenc}
    \usepackage{lmodern}
\fi

\usepackage[scaled=0.95]{helvet}
\renewcommand{\familydefault}{\sfdefault}

% Some settings:
\raggedright
\AtBeginEnvironment{adjustwidth}{\partopsep0pt}
\pagestyle{empty}
\setcounter{secnumdepth}{0}
\setlength{\parindent}{0pt}
\setlength{\topskip}{0pt}
\setlength{\columnsep}{0.15cm}
\pagenumbering{gobble}

\titleformat{\section}
    {\needspace{4\baselineskip}\bfseries\large}
    {}
    {0pt}
    {}
    [\vspace{1pt}\titlerule]

\titlespacing{\section}{
    -1pt
}{
    0.3 cm
}{
    0.2 cm
}

\renewcommand\labelitemi{$\vcenter{\hbox{\small$\bullet$}}$}

\newenvironment{highlights}{
    \begin{itemize}[
        topsep=0.10 cm,
        parsep=0pt,
        partopsep=0pt,
        itemsep=0pt,
        leftmargin=0 cm + 10pt
    ]
}{
    \end{itemize}
}

\newenvironment{highlightsforbulletentries}{
    \begin{itemize}[
        topsep=0.10 cm,
        parsep=0pt,
        partopsep=0pt,
        itemsep=0pt,
        leftmargin=10pt
    ]
}{
    \end{itemize}
}

\newcommand{\nobullet}{\item[]}

\newenvironment{onecolentry}{
    \begin{adjustwidth}{
        0 cm + 0.00001 cm
    }{
        0 cm + 0.00001 cm
    }
}{
    \end{adjustwidth}
}

\newenvironment{twocolentry}[2][]{
    \onecolentry
    \def\secondColumn{#2}
    \setcolumnwidth{\fill, 4.5 cm}
    \begin{paracol}{2}
}{
    \switchcolumn \raggedleft \secondColumn
    \end{paracol}
    \endonecolentry
}

\newenvironment{threecolentry}[3][]{
    \onecolentry
    \def\thirdColumn{#3}
    \setcolumnwidth{, \fill, 4.5 cm}
    \begin{paracol}{3}
    {\raggedright #2} \switchcolumn
}{
    \switchcolumn \raggedleft \thirdColumn
    \endonecolentry
    \end{paracol}
}

\newenvironment{header}{
    \setlength{\topsep}{0pt}\par\kern\topsep\centering\linespread{1.5}
}{
    \par\kern\topsep
}

\newcommand{\projectheader}[2]{%
    \href{#1}{\textbf{#2} \small\faExternalLink*}%
}

\begin{document}

\newcommand{\AND}{\unskip
    \cleaders\copy\ANDbox\hskip\wd\ANDbox
    \ignorespaces
}

\newsavebox\ANDbox
\sbox\ANDbox{$|$}

% =========================
% HEADER
% =========================

\begin{header}
    {\fontsize{25 pt}{25 pt}\selectfont Reece Palmer}\\[0.4cm]
    \normalsize
    \mbox{\href{mailto:reecelpalmer@gmail.com}{reecelpalmer@gmail.com}}%
    \kern 5.0 pt%
    \AND%
    \kern 5.0 pt%
    \mbox{\href{https://www.linkedin.com/in/reecepalmer}{linkedin.com/in/reecepalmer}}%
    \kern 5.0 pt%
    \AND%
    \kern 5.0 pt%
    \mbox{\href{https://github.com/rpalmr}{github.com/rpalmr}}%
\end{header}

% =========================
% EXPERIENCE
% (XYZ: accomplished [X], as measured by [Y], by doing [Z].
%  Anything in [square brackets] needs a real figure from you.)
% =========================

\section{Experience}

\begin{twocolentry}{
    Mar 2026 -- Present
}
    \textbf{Machine Learning Engineer}, AppliedCL
\end{twocolentry}

\vspace{0.10 cm}

\begin{onecolentry}
    \begin{highlights}
        \item Designed and implemented a protein analysis workflow to enable structure-informed searches with Foldseek, integrating pretrained AI models through PyTorch, Hugging Face Transformers, and Biopython to achieve 62.48\% accuracy in predicting per-residue structural features.
        \item Developed and evaluated AI models to identify interactions between proteins and protein complexes, systematically refining model architecture and training methods to achieve 93.56\% exploratory test AUROC, a measure of how well models distinguish interacting from non-interacting proteins.
        \item Advanced research into sequence-based interaction prediction by initiating ESM-2 fine-tuning on experimentally supported biological data, implementing reliable checkpoint recovery to support continued training across computing sessions.
        \item Established a repeatable evaluation framework using scikit-learn and MMseqs2, combining complementary AI representations to achieve 68.41\% mean validation AUROC and using prediction errors and protein similarity checks to guide subsequent experiments.
    \end{highlights}
\end{onecolentry}

\vspace{0.2 cm}

\begin{twocolentry}{
    Oct 2025 -- Sept 2026
}
    \textbf{Student Researcher}, University of Birmingham
\end{twocolentry}

\vspace{0.10 cm}

\begin{onecolentry}
    \begin{highlights}
        \item Improved Parkinson's detection from clinical voice data by [X\% / metric], by applying causal machine learning (de-confounding and transport methods).
        \item Measured how well models generalise across [N] patient groups, by building and evaluating Python ML pipelines for [N] model types.
    \end{highlights}
\end{onecolentry}

\vspace{0.2 cm}

\begin{twocolentry}{
    Sept 2024 -- Sept 2025
}
    \textbf{Independent Development}
\end{twocolentry}

\vspace{0.10 cm}

\begin{onecolentry}
    \begin{highlights}
        \item Entrepreneur's First Founder in Residence, Summer 2025 cohort: rapidly investigated and invalidated multiple ideas, learning quickly and building relationships with talented founders, before leaving the cohort without a business or co-founder to start my master's.
        \item Delivered 3 prototypes at London hackathons (a16z HACK UK, Imperial's IC HACK and Encode) within [24--48]-hour deadlines, by building [what] using [tech stack].
        \item Volunteered short time at different friend's London based startups usually following angel investments.
    \end{highlights}
\end{onecolentry}

\vspace{0.2 cm}


\begin{twocolentry}{
    Nov 2022 -- Oct 2023
}
    \textbf{Co-Founder}, \href{https://gininow.com}{GiniNow}
\end{twocolentry}

\vspace{0.10 cm}

\begin{onecolentry}
    \begin{highlights}
        \item Secured [N] partnerships and expert relationships for GINI, an online marketplace for hiring experts through video and text, by [outreach approach], as a 25\% equity co-founder.
        \item Kept a team of 7 delivering, with [N] months of runway managed on a [amount] budget, by running day-to-day operations and company finances.
        \item Raised a [amount] pre-seed round to fund platform development, by pitching to [N] investors.
    \end{highlights}
\end{onecolentry}

\vspace{0.2 cm}


% =========================
% VOLUNTEERING & TRUSTEE EXPERIENCE
% =========================

\section{Volunteering \& Trustee Experience}

\begin{twocolentry}{
    Month Year -- Present
}
    \textbf{Trustee}, Organisation Name
\end{twocolentry}

\vspace{0.10 cm}

\begin{onecolentry}
    \begin{highlights}
        \item Helped [outcome, e.g. grow funding / improve governance] at a [charity type] supporting [beneficiaries], measured by [figure], by [action, e.g. leading a review or advising on data and technology].
        \item Oversaw [budget / programme / decision] worth [amount or scale], by [committee role or responsibility].
    \end{highlights}
\end{onecolentry}

\vspace{0.2 cm}

\begin{twocolentry}{
    Month Year -- Month Year
}
    \textbf{Volunteer Role}, Organisation Name
\end{twocolentry}

\vspace{0.10 cm}

\begin{onecolentry}
    \begin{highlights}
        \item Achieved [outcome], measured by [figure, e.g. people reached, hours, funds raised], by [what you did].
    \end{highlights}
\end{onecolentry}


% =========================
% EDUCATION
% =========================

\section{Education}

\begin{twocolentry}{
    Sept 2025 -- Sept 2026
}
    \textbf{MSc Computer Science}, University of Birmingham
\end{twocolentry}

\vspace{0.10 cm}

\begin{onecolentry}
    \begin{highlights}
        \item \textbf{Overall Grade: 71\%}
        \item \textbf{Dissertation (75\%)}: Causal Deconfounding and Transport for Parkinson's Voice Detection (Supervised by Dr. Max Little). Improved Parkinson's detection across patient populations by [X\% / metric], by using causal machine learning to reduce confounding in voice data.
        \item \textbf{Key Modules:} Machine Learning, Artificial Intelligence, Computer Systems, Data Structures \& Algorithms.
    \end{highlights}
\end{onecolentry}

\vspace{0.2 cm}

\begin{twocolentry}{
    Sept 2020 -- July 2023
}
    \textbf{BA Business}, University of Portsmouth
\end{twocolentry}

\vspace{0.10 cm}

\begin{onecolentry}
    \begin{highlights}
        \item \textbf{Overall Grade: 68\%}
        \item \textbf{Dissertation (70\%)}: Strategy, Enterprise and Innovation
        \item \textbf{Key Modules:} Business Analysis \& Decision Making, Strategic Management, Accounting
    \end{highlights}
\end{onecolentry}

% =========================
% PROJECTS
% =========================


\end{document}